\providecommand{\itescastudentname}{Martin Jonathan de la Cruz}
\providecommand{\itescastudentshort}{M. J. de la Cruz}
\providecommand{\itescastudentid}{ITESCA-MGA}
\providecommand{\itescafaculty}{Instituto Tecnologico Superior de Cajeme}
\providecommand{\itescafacultyname}{Maestria en Gestion Administrativa}
\providecommand{\itescacoursename}{Maestria en Gestion Administrativa}
\providecommand{\itescacoursecode}{MGA}
\providecommand{\itescadepartment}{Trayectoria academica por carrera}
\providecommand{\itescateachingfigure}{Coordinacion academica por definir}
\providecommand{\itescaprogramlevel}{Maestria}
\providecommand{\itescaprogramperiod}{Periodo academico por definir}
\providecommand{\itescamodality}{Modalidad por definir}
\providecommand{\itescaprofessionalprofile}{La Maestria en Gestion Administrativa exige analizar decisiones organizacionales con criterio, evidencia, pertinencia institucional y enfoque de mejora.}
\providecommand{\itescatransfercontext}{Diagnostico organizacional, planeacion, gestion de procesos y toma de decisiones administrativas}
\providecommand{\itescalocation}{Cajeme, Sonora}
\providecommand{\itescalocationname}{Cajeme, Sonora}
\providecommand{\itescabibliography}{bibliografia-itesca-mga}
\providecommand{\itescageneralbib}{bibliografia-itesca-mga.bib}
\providecommand{\itescacoursebib}{bibliografia-itesca-mga.bib y referencias-itesca-mga/}
\def\itescacoursename {Primer ingreso}
\def\itescacoursecode {MGA-PI}
\def\itescafaculty {Maestria en Gestion Administrativa}
\def\itescafacultyname {Maestria en Gestion Administrativa}
\def\itescadepartment {Primer ingreso}
\def\itescateachingfigure {Figura docente por definir}
\def\itescaprofessionalprofile {En la Maestria en Gestion Administrativa, el primer ingreso debe convertir la induccion en habitos de analisis, planeacion, evidencia y mejora organizacional.}
\def\itescatransfercontext {Habitos de posgrado, diagnostico administrativo, comunicacion ejecutiva y documentacion de decisiones}
\def\itescabibliography {ITESCA/maestria-en-gestion-administrativa/primer-ingreso-mga/primer-ingreso-mga}
\def\itescacoursebib {primer-ingreso-mga.bib}

\def\itescacoursename {Primer ingreso}
\def\itescacoursecode {MGA-PI}
\def\itescafaculty {Maestria en Gestion Administrativa}
\def\itescafacultyname {Maestria en Gestion Administrativa}
\def\itescadepartment {Primer ingreso}
\def\itescateachingfigure {Figura docente por definir}
\def\itescaprofessionalprofile {En la Maestria en Gestion Administrativa, el primer ingreso debe convertir la induccion en habitos de analisis, planeacion, evidencia y mejora organizacional.}
\def\itescatransfercontext {Habitos de posgrado, diagnostico administrativo, comunicacion ejecutiva y documentacion de decisiones}
\def\itescabibliography {ITESCA/maestria-en-gestion-administrativa/primer-ingreso-mga/primer-ingreso-mga}
\def\itescacoursebib {primer-ingreso-mga.bib}

% ==============================================================================
% PLANTILLA MAESTRA GENERAL - ITESCA
% ============================================================================== 
% Uso:
% 1. Duplica el archivo contenedor de la materia o actividad y actualiza sus
%    metadatos editables.
% 2. Conserva el enfoque de trabajo: no pegar instrucciones completas dentro del
%    documento final; usar la planeacion como lista de cumplimiento.
% 3. Construye productos visuales en LaTeX/TikZ cuando la actividad lo pida.
% 4. Mantiene citas autor-anio con natbib y bibliografia en BibTeX.
% 5. Cierra con criterio propio, no con un resumen plano.
%
% Formato institucional recomendado para ITESCA:
% - Fuente sans serif tipo Helvetica.
% - Interlineado 1.5.
% - Portada institucional, abstract, desarrollo, postura y conclusion.
% - Citas autor-anio con natbib: usar \citep{clave} o \citet{clave}.
% - Bibliografia en BibTeX, consistente con el producto real.
% - Comentarios operativos dentro del archivo para sostener la redaccion.
% - Identidad institucional visible mediante logotipo oficial, marca de agua y
%   vinculacion con la carrera o el contexto profesional.
%
% Criterios generales de cumplimiento:
% - Responder exactamente a la actividad solicitada por la materia o docente.
% - Convertir la consigna en problema, analisis, evidencia y consecuencia.
% - Relacionar el tema con la carrera, la materia y el contexto institucional.
% - Analizar, interpretar y tomar postura; no solo recopilar definiciones.
% - Mostrar producto visible cuando la actividad pida cuadro, esquema, tabla,
%   mapa, diagrama, evidencia o procedimiento.
% - Cerrar con criterio propio, aprendizaje y aplicacion transferible.
%
% Regla para productos visuales:
% - Si la actividad pide cuadro comparativo, mapa, linea de tiempo, matriz,
%   tabla, diagrama o evidencia analitica, construirlo preferentemente en
%   LaTeX/TikZ para mantener nitidez, control y edicion futura.
% - Si la actividad exige herramienta externa, insertar una imagen nitida y
%   conservar la explicacion dentro del documento.
%
% Criterios de redaccion personal:
% - No escribir para llenar espacio; escribir para sostener una idea.
% - Convertir informacion en criterio: que significa, como funciona, que limite
%   presenta, que riesgo aparece y que consecuencia produce.
% - Formula util: concepto -> funcion -> problema -> consecuencia -> postura.
% - La postura personal no debe ser opinion suelta. Debe tener posicion, razon y efecto.
% - La conclusion no debe repetir el desarrollo: debe sintetizar y proyectar.
%
% Control editorial ITESCA:
% - Cada actividad debe tener una ficha de lectura: consigna, producto,
%   evidencia, criterio de evaluacion y transferencia profesional.
% - La identidad institucional no debe limitarse al logotipo. Debe aparecer en
%   la forma de argumentar: precision tecnica, pertinencia academica,
%   vinculacion con carrera y utilidad profesional.
% - Evitar la respuesta generica. El documento debe mostrar que problema atiende
%   la actividad, que decision metodologica se tomo y que evidencia permite
%   sostener el resultado.
% - Cuando la actividad sea de induccion, diagnostico o primer ingreso, traducir
%   el contenido a habitos universitarios: lectura, organizacion, evidencia,
%   integridad academica, autonomia y seguimiento institucional.
% - Cuando la actividad sea disciplinar, traducir el contenido a un proceso,
%   sistema, caso, herramienta, decision administrativa, tecnica o profesional.
%
% Regla de nivel minimo:
% - Si la actividad solo pide identificar o describir, elevar el producto a
%   comprension aplicada: explicar funcion, limite y ejemplo.
% - Si pide cuadro, mapa, matriz, evidencia o conclusion, llegar por lo menos a
%   analisis: comparar criterios, justificar diferencias y cerrar con postura.
% - Si pide propuesta, procedimiento o solucion, llegar a evaluacion/creacion:
%   definir criterios, construir producto propio y justificar su pertinencia.
%
% Voz editorial esperada:
% - Academica, directa, institucional y funcional.
% - No sonar como copia de fuentes ni como resumen escolar.
% - Usar criterio personal con moderacion: posicion, razon y efecto.
% - Antes de redactar, decidir la funcion de cada seccion: abrir problema,
%   delimitar criterio, demostrar con evidencia, valorar o concluir.
%
% Nivel cognitivo esperado:
% - Comprender: explicar el tema con lenguaje propio.
% - Aplicar: llevar el concepto a un caso, proceso, procedimiento o producto.
% - Analizar: distinguir componentes, tensiones, limites y decisiones.
% - Evaluar: sostener un criterio con razones y consecuencias.
% - Crear: construir un producto propio cuando la actividad lo solicite.
%
% Checklist antes de entregar:
% [ ] Titulo, actividad, asignatura, codigo, docente y fecha actualizados.
% [ ] El objetivo coincide con la consigna real.
% [ ] El producto solicitado aparece visible y completo.
% [ ] La introduccion plantea un problema y no solo anuncia el tema.
% [ ] El desarrollo explica, conecta y analiza.
% [ ] El texto aterriza el tema a la carrera o al contexto profesional.
% [ ] Hay citas autor-anio y bibliografia consistente.
% [ ] La conclusion responde que se comprendio y que criterio se sostiene.
% [ ] Quitar o comentar \pendiente cuando el documento final este terminado.
% ============================================================================== 

% ------------------------------------------------------------------------------
% METADATOS EDITABLES
% ------------------------------------------------------------------------------
\providecommand{\itescadocumenttitle}{Actividad 2: Cuestionario}
\providecommand{\itescadocumentsubtitle}{Respuestas justificadas}
\providecommand{\itescadocumentsubject}{Cuestionario academico}
\providecommand{\itescastudentname}{Martin Jonathan de la Cruz}
\providecommand{\itescastudentid}{ITESCA}
\providecommand{\itescacoursename}{Nombre de la asignatura}
\providecommand{\itescacoursecode}{ITESCA}
\providecommand{\itescafaculty}{Programa academico ITESCA}
\providecommand{\itescadepartment}{Departamento academico}
\providecommand{\itescateachingfigure}{Nombre por definir}
\providecommand{\itescastudentemail}{Dato no proporcionado}
\providecommand{\itescaprogramlevel}{Maestria}
\providecommand{\itescaprogramperiod}{Primer ingreso 2026}
\providecommand{\itescamodality}{Posgrado}
\providecommand{\itescaactivityname}{Actividad 2. Cuestionario resuelto}
\providecommand{\itescasessionname}{Primer ingreso}
\providecommand{\itescablockname}{Fundamentos administrativos y comerciales}
\providecommand{\itescamainproduct}{Cuestionario resuelto con justificacion conceptual y numerica}
\providecommand{\itescaprofessionalprofile}{En la Maestria en Gestion Administrativa, el primer ingreso debe convertir la induccion en habitos de analisis, planeacion, evidencia y mejora organizacional.}
\providecommand{\itescainstitutionalaxis}{Formacion academica, identidad institucional y transferencia profesional}
\providecommand{\itescaevidencefocus}{Evidencia academica, producto visible y criterio propio}
\providecommand{\itescaqualitycriterion}{Precision conceptual, procedimiento correcto, evidencia y cierre analitico}
\providecommand{\itescatransfercontext}{Habitos de posgrado, diagnostico administrativo, comunicacion ejecutiva y documentacion de decisiones}
\providecommand{\itescalocation}{Cajeme, Sonora}
\providecommand{\itescabibliography}{ITESCA/maestria-en-gestion-administrativa/primer-ingreso-mga/primer-ingreso-mga}
\providecommand{\itescaasset}{ITESCA/assets-itesca/itesca-monograma.png}
\providecommand{\itescaofficialasset}{ITESCA/assets-itesca/web/logo-itesca-contorno.png}
\providecommand{\itescasecondaryasset}{ITESCA/assets-itesca/itesca-monograma.png}
\providecommand{\itescacampusurl}{https://www.itesca.edu.mx/}
\providecommand{\itescaportalurl}{https://www.itesca.edu.mx/portalacademico/}
\providecommand{\itescavirtualurl}{https://www.itesca.edu.mx/modulos/oferta/ead/index.php}

% CREACION DEL DOCUMENTO
\documentclass[
	spanish,
	letterpaper, oneside
]{article}

% INFORMACION DEL DOCUMENTO
\def\documenttitle {\itescadocumenttitle}
\def\documentsubtitle {\itescadocumentsubtitle}
\def\documentsubject {\itescadocumentsubject}

\def\documentauthor {\itescastudentname}
\def\coursename {\itescacoursename}
\def\coursecode {\itescacoursecode}

\def\universityname {Instituto Tecnologico Superior de Cajeme}
\def\universityfaculty {\itescafaculty}
\def\universitydepartment {\itescadepartment}
\def\universitydepartmentimage {\itescaofficialasset}
\def\universitydepartmentimagecfg {width=5.2cm}
\def\universitylocation {\itescalocation}

% INTEGRANTES, PROFESORES Y FECHAS
\def\authortable {
	\begin{tabular}{ll}
		\textbf{Alumno:}
		& \begin{tabular}[t]{l}
			 \itescastudentname \\
		\end{tabular} \\
		Matricula: & \itescastudentid \\
		Correo institucional: & \itescastudentemail \\
		Asignatura: & \itescacoursename \\
		Codigo: & \itescacoursecode \\
		Nivel: & \itescaprogramlevel \\
		Periodo: & \itescaprogramperiod \\
		Modalidad: & \itescamodality \\

		\textit{Figura docente:}
		& \begin{tabular}[t]{l}
			\itescateachingfigure
		\end{tabular} \\
		Sesion/Bloque: & \itescasessionname\ / \itescablockname \\
		Producto: & \itescamainproduct \\
		& \\
		\multicolumn{2}{l}{Fecha de realizacion: \today} \\
		\multicolumn{2}{l}{\universitylocation}
	\end{tabular}
}

% IMPORTACION DEL TEMPLATE BASE
\input{template}

% FORMATO GENERAL
\usepackage{helvet}
\renewcommand{\familydefault}{\sfdefault}

% IDENTIDAD VISUAL ITESCA
\definecolor{itescaRedDark}{HTML}{8F1117}
\definecolor{itescaRed}{HTML}{D70F18}
\definecolor{itescaCharcoal}{HTML}{141414}
\definecolor{itescaSilver}{HTML}{D9D9D9}
\definecolor{itescaPaper}{HTML}{F7F7F7}
\definecolor{itescaInk}{HTML}{181818}
\definecolor{itescaMuted}{HTML}{6A6A6A}
\colorlet{itescaGreenDark}{itescaRedDark}
\colorlet{itescaGreen}{itescaCharcoal}
\colorlet{itescaGold}{itescaSilver}
\def\portraittitlecolor {itescaRedDark}

% TABLAS Y PRODUCTOS VISUALES
\usepackage{array}
\usepackage{longtable}
\usepackage{pdflscape}
\usepackage{tikz}
\newcolumntype{L}[1]{>{\raggedright\arraybackslash}p{#1}}
\usetikzlibrary{
	arrows.meta,
	positioning,
	calc,
	matrix,
	fit,
	shapes.geometric,
	shadows.blur
}

% CITAS EN FORMATO AUTOR-ANIO, COMPATIBLE CON APA 7 EN EL CUERPO DEL TEXTO
\setcitestyle{authoryear,open={(},close={)}}

% MARCA DE AGUA EN PORTADA
\def\coverwatermarkenabled {true}
\def\coverwatermarkimage {\itescaasset}
\def\coverwatermarkopacity {0.13}
\def\coverwatermarkwidth {0.58\paperwidth}
\def\coverwatermarkxshift {0cm}
\def\coverwatermarkyshift {-0.10cm}

\newcommand{\insertcoverwatermark}{%
	\ifthenelse{\equal{\coverwatermarkenabled}{true}}{%
		\AddToShipoutPictureBG*{%
			\begin{tikzpicture}[remember picture,overlay]
				\node[
					opacity=\coverwatermarkopacity,
					inner sep=0pt,
					xshift=\coverwatermarkxshift,
					yshift=\coverwatermarkyshift
				] at (current page.center) {%
					\includegraphics[width=\coverwatermarkwidth]{\coverwatermarkimage}%
				};
			\end{tikzpicture}%
		}%
	}{}%
}

\newcommand{\insertcoverinstitutionalfooter}{%
	\AddToShipoutPictureFG*{%
		\begin{tikzpicture}[remember picture,overlay]
			\fill[itescaCharcoal] (current page.south west) rectangle ([yshift=1.35cm]current page.south east);
			\fill[itescaRed] ([yshift=1.35cm]current page.south west) rectangle ([yshift=1.47cm]current page.south east);
			\node[anchor=south west] at ([xshift=0.95cm,yshift=0.16cm]current page.south west) {%
				\includegraphics[height=0.88cm]{\itescasecondaryasset}%
			};
			\node[
				anchor=south east,
				text=white,
				align=right,
				font=\sffamily\bfseries\small
			] at ([xshift=-0.95cm,yshift=0.28cm]current page.south east) {%
				Instituto Tecnologico Superior de Cajeme\\
				\textcolor{itescaSilver}{\itescalocation \textperiodcentered\ \href{\itescacampusurl}{itesca.edu.mx}}%
			};
		\end{tikzpicture}%
	}%
}

\begin{document}

% PORTADA
\insertcoverwatermark
\insertcoverinstitutionalfooter
\templatePortrait

% CONFIGURACION DE PAGINA Y ENCABEZADOS
\templatePagecfg
\onehalfspacing

% RESUMEN / ABSTRACT
\begin{abstractd}
	\textbf{Objetivo:} Analizar y resolver un cuestionario de fundamentos de administracion, finanzas, mercadotecnia y proyectos de inversion para demostrar comprension aplicada de conceptos esenciales en la toma de decisiones organizacionales.

	\textbf{Producto elaborado:} Cuestionario resuelto con respuestas justificadas y procedimientos sinteticos en los reactivos numericos.

	\textbf{Enfoque de analisis:} El documento interpreta la actividad como un ejercicio de integracion entre proceso administrativo, control financiero, logica de mercado y criterio operativo, no como una acumulacion mecanica de respuestas.

	\textbf{Vinculacion institucional:} La actividad fortalece habitos de posgrado vinculados con \itescafaculty, ya que exige argumentar, verificar resultados y relacionar conceptos con \itescatransfercontext.

	\textbf{Nivel cognitivo:} Comprender, aplicar y analizar.
\end{abstractd}

% CONFIGURACIONES FINALES
\templateFinalcfg

% ======================= INICIO DEL DOCUMENTO =======================

\section{Introduccion}

Responder un cuestionario de administracion no consiste en repetir definiciones, sino en demostrar que los conceptos sirven para interpretar decisiones reales dentro de una organizacion. Cuando una actividad integra proceso administrativo, razon financiera, mercadotecnia y evaluacion de proyectos, lo que en realidad pone a prueba es la capacidad para conectar funciones que en la practica operan como un solo sistema. Desde esta perspectiva, la administracion organiza recursos y procesos, mientras que la mercadotecnia orienta el valor hacia el mercado y las finanzas verifican la viabilidad y el control de los resultados \citep{amelicaProcesoAdmin2019,snhuAdminMerca2023}.

En el contexto de primer ingreso de la Maestria en Gestion Administrativa, este tipo de ejercicio cumple una funcion formativa clara: obligar a pasar de la memorizacion a la comprension aplicada. Cada reactivo exige identificar relaciones entre planeacion, direccion, segmentacion, rentabilidad, flujo de efectivo y proyectos de inversion, es decir, entre decisiones que afectan directamente la competitividad de las organizaciones \citep{amelicaProcesoAdmin2019,unadmFundAdmin}.

El criterio que guia este documento sostiene que una respuesta correcta solo tiene valor academico cuando puede justificarse con logica administrativa y no solo con reconocimiento superficial del termino. Por ello, el cuestionario se presenta resuelto, pero tambien interpretado, de modo que el producto final muestre conocimiento operativo y no un listado desarticulado de opciones.

\section{Cuestionario resuelto}

El presente producto corresponde a un cuestionario resuelto. Cada reactivo incluye el enunciado textual, la respuesta correcta con letra y redaccion completa, asi como una justificacion breve. En los casos numericos se agrega el procedimiento esencial para mostrar de donde sale el resultado.

\begin{center}
\textbf{\large Actividad 2. Cuestionario resuelto}
\end{center}

\clearpage
\pagestyle{empty}
\thispagestyle{empty}
\begin{landscape}
\begingroup
\fontsize{6.6}{7.5}\selectfont
\setlength{\tabcolsep}{3pt}
\renewcommand{\arraystretch}{1.18}
\begin{longtable}{@{}L{0.03\linewidth}L{0.41\linewidth}L{0.24\linewidth}L{0.26\linewidth}@{}}
\caption{Actividad 2. Cuestionario resuelto}\label{tab:cuestionario-actividad-2}\\
\toprule
\textbf{No.} & \textbf{Pregunta} & \textbf{Respuesta} & \textbf{Justificacion o procedimiento} \\
\midrule
\endfirsthead
\toprule
\textbf{No.} & \textbf{Pregunta} & \textbf{Respuesta} & \textbf{Justificacion o procedimiento} \\
\midrule
\endhead
\midrule
\multicolumn{4}{r}{Continua en la siguiente pagina} \\
\endfoot
\bottomrule
\endlastfoot
1 & Son aquellas empresas que transforman las materias primas en productos terminados. & c) Manufactureras & Son manufactureras porque convierten insumos en bienes finales. \\
2 & Elementos que disputan las mismas entradas (proveedores) y las mismas salidas (clientes) de la organizacion. & c) Competidores & Son los competidores, ya que buscan los mismos proveedores y clientes. \\
3 & Son aquellas empresas que transforman las materias primas en productos terminados. & c) Manufactureras & Se repite el concepto de empresa manufacturera. \\
4 & Es la secuencia de pasos que permiten llevar a cabo la administracion de una organizacion para el logro de sus objetivos. & c) Proceso administrativo & Es el proceso administrativo: planeacion, organizacion, direccion y control. \\
5 & Busca la ejecucion de las actividades encaminadas al logro de los objetivos, mediante el esfuerzo del recursos humanos y la eficiencia de sus recursos materiales, financieros y tecnologicos. & c) Administracion & Esa definicion corresponde a la administracion. \\
6 & En esta fase del proceso administrativo se agudiza la toma de decisiones, se ve todo lo relacionado con la integracion, motivacion, comunicacion y supervision de los integrantes de la organizacion con el proposito de contribuir al logro de los objetivos. & d) Direccion & La direccion coordina y conduce al personal hacia los objetivos. \\
7 & Interrogante de la organizacion que sirve para determinar los metodos de actuacion. & c) \textquestiondown Como? & La pregunta \emph{como} define la manera de ejecutar. \\
8 & Tipo de planeacion que se aplica en los niveles inferiores de la organizacion y consiste en la planeacion de las actividades cotidianas. & d) Planeacion operativa & La planeacion operativa organiza las tareas diarias. \\
9 & Tipo de planeacion que requiere del establecimiento de planes maestros para definir el destino de una empresa, actividad que recae en el mas alto nivel de la organizacion. & a) Planeacion estrategica & La planeacion estrategica fija objetivos generales y largo plazo. \\
10 & Es el elemento del proceso administrativo que vigila el rumbo a donde se encamina la organizacion por medio de la autoridad, el liderazgo, comunicacion, asi como los cambios organizacionales e individuales. & d) Direccion & La direccion orienta la accion organizacional. \\
11 & Elemento del proceso administrativo que consiste en verificar si todo ocurre de conformidad con el plan adoptado, con las instrucciones emitidas y con los principios establecidos. Tiene como fin senalar las debilidades y errores a fin de rectificarlos e impedir que se produzcan nuevamente. & a) Control & El control compara resultados y corrige desviaciones. \\
12 & Son aquellas empresas que efectuan ventas en gran escala a otras empresas tanto al menudeo como al detalle. & b) Comerciales-mayoristas & Son comerciales mayoristas por su volumen de venta. \\
13 & Selecciona el orden de los pasos de la planeacion. & c) 1. Deteccion de la oportunidad. 2. Establecer objetivos. 3. Desarrollar premisas. 4. Determinacion de los cursos alternativos de accion. 5. Evaluacion de los cursos alternativos de accion. 6. Seleccionar un curso. 7. Formular planes derivados. 8. Cuantificar planes mediante presupuestos. & Resume la secuencia clasica del proceso de planeacion. \\
14 & Es uno de los criterios mas utilizados para clasificar a las empresas, el que de acuerdo al tamano de la misma se establece que puede ser pequena, mediana o grande. & a) Clasificacion por su magnitud & Se clasifica por su magnitud. \\
15 & Clasificacion que se basa en el origen de las aportaciones y del caracter a quien se dirijan sus actividades las empresas. & b) Clasificacion por el origen de su capital & Corresponde al origen del capital. \\
16 & Factor del medio ambiente que contiene los elementos relacionados con los valores y costumbres que prevalecen en un grupo o region. & c) Factores culturales & Son factores culturales. \\
17 & Son parte del ambiente general que incluye la inflacion, politica fiscal de gastos de gobierno, tamano del mercado, poder de compra, entre otros. & c) Factores economicos & Son factores economicos. \\
18 & Es el area que se encarga del optimo control, manejo de recursos economicos de la empresa, incluye la obtencion de recursos financieros tanto internos como externos, necesarios para alcanzar los objetivos y metas empresariales. & c) Finanzas & Finanzas obtiene y administra recursos monetarios. \\
19 & La empresa Fashion, S. A. de C. V., se dedica a la produccion y venta de pantalones de mezclilla para el mercado femenino; la matriz se encuentra ubicada en la Ciudad de Guadalajara, Jalisco, y cuenta ademas con 15 sucursales en el resto de la Republica Mexicana. Una de las sucursales de la empresa proporciona la informacion correspondiente para el proximo trimestre: $60{,}000.00$ por concepto de sueldos a los directores de la planta fabril, $80.00$ por unidad en materia prima, $95{,}000.00$ de sueldos del area administrativa, precio de venta de $200.00$ por prenda, prestamo de $77{,}500.00$ para cubrir su nomina con tasa anual de 23\% y volumen estimado de ventas de 1,500 unidades. Se pide determinar el importe de los intereses generados por el prestamo y el margen de utilidad neta. & c) a) $4{,}456.20$; b) $6.84\%$ & Interes trimestral $=77{,}500(0.23/4)=4{,}456.25$; utilidad neta $=300{,}000-120{,}000-60{,}000-95{,}000-4{,}456.25=20{,}543.75$; margen $=20{,}543.75/300{,}000\times 100\approx 6.84\%$. \\
20 & El corporativo La Gaseosa tiene bajo su control financiero las siguientes 3 empresas con sus respectivas utilidades: Vidrios el Cristal $3{'}000{,}000$ anuales, Refresquera La Naranjita $8{'}000{,}000$ anuales y Galletas La Antiguita $6{'}000{,}000$ anuales. Si el corporativo presenta un margen de utilidad neta de 80\%, \textquestiondown a cuanto ascienden las ventas anuales de las empresas? & c) $21{,}250{,}000.00$ & Utilidad total $=3{,}000{,}000+8{,}000{,}000+6{,}000{,}000=17{,}000{,}000$; ventas $=17{,}000{,}000/0.80=21{,}250{,}000$. \\
21 & Determina el costo unitario del producto si se tiene la siguiente informacion: la materia prima requiere de 5 unidades por producto y el costo de la materia prima es de $35.00$ cada unidad; requiere de 5 horas de mano de obra para elaborar el producto y la hora se paga a $55.00$; los gastos indirectos de fabricacion ascienden a $500{,}000.00$ y la aplicacion a cada producto debe calcularse de acuerdo al total de horas de mano de obra; el total de horas en la produccion fue de 250,000 horas. & b) $460.00$ & MP $=5(35)=175$; MOD $=5(55)=275$; GIF por hora $=500{,}000/250{,}000=2$; GIF unitario $=5(2)=10$; total $=460$. \\
22 & Representa las deudas que tiene el negocio con otras empresas y las obligaciones contraidas. & a) Pasivo & Corresponde al pasivo. \\
23 & Representa el derecho que los socios tienen sobre la empresa asi como las aportaciones de los duenos. & d) Capital & Representa el capital contable. \\
24 & Representan los bienes y derechos propiedad de una empresa, mismos que se tienen en dinero, se tienen para vender y para trabajar. & c) Activo & Esa definicion corresponde al activo. \\
25 & Los cuatro estados financieros clave que deben reportarse a los accionistas son: 1) el estado de perdidas y ganancias o estado de resultados, 2) el balance general [estado de situacion financiera], 3) el estado de patrimonio de los accionistas y 4) \rule{1cm}{0.15mm}. & c) El estado de flujos de efectivo & Es el estado de flujos de efectivo. \\
26 & Tomando como base las siguientes cuentas del Balance General de TUEMPRESA, calcula la razon financiera prueba del acido: Efectivo $15{,}000$, Clientes $59{,}400$, Inventario $65{,}000$, Planta y equipo $172{,}000$, Equipo y transporte $311{,}000$, Proveedores $8{,}500$, Documentos por pagar $17{,}700$, IVA por pagar $3{,}168$. & c) 2.53 & $(15{,}000+59{,}400)/(8{,}500+17{,}700+3{,}168)=74{,}400/29{,}368\approx 2.53$. \\
27 & Elija la opcion que contiene los enunciados que incluyen las caracteristicas especificas de los flujos de efectivo: 1. La depreciacion se calcula por medio del codigo fiscal. 2. El estado de flujos de efectivo se divide en flujos operativos, de inversion y de financiamiento. 3. Los estados pro forma se usan para pronosticar y analizar el nivel de rentabilidad. 4. Desde un punto de vista estrictamente financiero, el flujo de efectivo operativo de una empresa excluye los intereses. 5. El flujo de efectivo libre de una empresa es el efectivo disponible para los acreedores y propietarios. & b) 3, 4 y 5 & El reactivo tiene inconsistencia; conceptualmente la combinacion correcta seria 2, 4 y 5. Entre las opciones dadas, la mas cercana es B porque conserva 4 y 5. \\
28 & \textquestiondown Cuales son los componentes mas comunes de las entradas de efectivo de una empresa durante un periodo financiero especifico? & c) Ventas en efectivo y cuentas por cobrar & Provienen de ventas en efectivo y recuperacion de cuentas por cobrar. \\
29 & Alejandro deposita $18{,}000$ en una cuenta de ahorros que paga el 6\% de interes compuesto anual. Si desea dejar el dinero durante 8 anos, \textquestiondown cuanto dinero tendra en la cuenta al termino del periodo? & d) $28{,}689$ & $VF=18{,}000(1.06)^8\approx 28{,}689$. \\
30 & Andrea desea calcular el valor presente de $150{,}000$ que recibira dentro de 12 anos. Suponiendo que su costo de oportunidad es del 4.5\%, \textquestiondown cual es el valor presente? & d) $88{,}450$ & $VP=150{,}000/(1.045)^{12}\approx 88{,}450$. \\
31 & Francisco desea hacer un deposito de $10{,}000$ y dejarlo durante 2 anos en un banco que le da el 8\% de interes capitalizable trimestralmente. \textquestiondown Cuanto dinero tendra Francisco al vencimiento? & c) $11{,}717$ & $VF=10{,}000(1+0.08/4)^8\approx 11{,}717$. \\
32 & La tabla muestra los costos y ponderaciones que tiene TUEMPRESA para los diversos tipos de capital: costo de deuda 5.6\% con porcentaje 40\%; costo de acciones preferentes 10.6\% con porcentaje 10\%; costo de ganancia retenidas 13.0\% con porcentaje 50\%. Con estos datos diga \textquestiondown cual es el costo de capital promedio ponderado (CCPP)? & b) 9.8\% & $CCPP=(5.6\%)(0.40)+(10.6\%)(0.10)+(13\%)(0.50)=9.8\%$. \\
33 & La tabla muestra los costos y ponderaciones que tiene TUEMPRESA para los diversos tipos de capital: costo de deuda 5.6\% con porcentaje 40\%; costo de acciones preferentes 10.6\% con porcentaje 10\%; costo de ganancia retenidas 13.0\% con porcentaje 50\%. Con estos datos diga \textquestiondown cual es el costo de capital promedio ponderado (CCPP)? & b) 9.8\% & Se repite el mismo calculo: $9.8\%$. \\
34 & Tomando como base las siguientes cuentas del Balance General de TUEMPRESA, calcula la razon financiera circulante: Efectivo $15{,}000$, Clientes $59{,}400$, Inventario $65{,}000$, Planta y equipo $172{,}000$, Equipo y transporte $311{,}000$, Proveedores $8{,}500$, Documentos por pagar $17{,}700$, IVA por pagar $3{,}168$. & a) 4.75 & $(15{,}000+59{,}400+65{,}000)/(8{,}500+17{,}700+3{,}168)=139{,}400/29{,}368\approx 4.75$. \\
35 & Segun esta teoria, las decisiones de compra son resultado de calculos economicos, racionales y conscientes. El comprador trata de gastar su dinero en mercancias que le proporcionen utilidad de acuerdo con sus gustos. & c) Modelo economico de Alfred Marshall & Se atribuye al modelo economico de Alfred Marshall. \\
36 & Su objetivo no es solo ser elegido de entre los competidores, sino ademas lograr cobrar mas por un producto similar; esto se logra mediante la creacion de valor en el costo o desempeno. & a) Diferenciacion & Esa estrategia es la diferenciacion. \\
37 & El proceso social y administrativo por el cual los grupos e individuos satisfacen sus necesidades al crear e intercambiar bienes y servicios. & a) Mercadotecnia & Es la definicion de mercadotecnia. \\
38 & Una buena segmentacion debe tener como resultado subgrupos o segmentos de mercado con las siguientes caracteristicas. & c) Todas las opciones & Debe ser homogenea internamente, heterogenea entre segmentos y operable. \\
39 & Estrategia mediante la cual se colocan los productos en tantos puntos de venta como sea posible. & d) Distribucion intensiva & Se llama distribucion intensiva. \\
40 & Es el conjunto de caracteristicas de un producto que determinan su capacidad de satisfacer necesidades. & d) Calidad & Esa capacidad se denomina calidad. \\
41 & Esta etapa de la vida del producto se caracteriza por pico de ventas, bajo coste por consumidor, altos beneficios; el objetivo se encuentra en la mayoria media y la competencia empieza a declinar. & d) Madurez & Corresponde a la madurez. \\
42 & Abarca una gran variedad de esfuerzos de comunicacion para contribuir a actitudes y opiniones generalmente favorables hacia una organizacion y sus productos. Los objetivos pueden ser clientes, accionistas, una organizacion gubernamental o un grupo de interes especial. & d) Relaciones Publicas & Son relaciones publicas. \\
43 & Es un tipo de estructura en donde encontramos a muchos compradores y vendedores que comercian con un producto homogeneo. & c) Competencia perfecta & Es competencia perfecta. \\
44 & Las cantidades de un producto que las personas estan dispuestas y en condiciones de adquirir a los diferentes precios en un momento determinado. & c) Demanda & Esa definicion corresponde a la demanda. \\
45 & \textquestiondown Cual de las siguientes afirmaciones es falsa? & c) Los estudios previos son suficientes para tomar la decision de continuar o abandonar el proyecto. & Es falsa porque los estudios previos no siempre bastan para decidir. \\
46 & \textquestiondown Quien realiza un proyecto? & a) Todas las opciones & Puede ser consultor, empresa de ingenieria u organizacion promotora. \\
47 & Son aquellos gastos que se hacen, generalmente al principio del proyecto; consiste en montar la empresa y dejarla lista para llevar a cabo las actividades productivas. & c) Inversion & Se trata de la inversion inicial para poner en marcha el proyecto. \\
48 & Se utiliza para evaluar inversiones de tipo social y los criterios. & d) Costo/Beneficio & El criterio costo-beneficio compara utilidad social y costo. \\
49 & Metodo de evaluacion que se utiliza para analizar el remplazo de equipo, tanto en el sector publico como en el privado. & d) Costo Anual & El costo anual facilita comparar alternativas de sustitucion. \\
50 & Es la tasa de crecimiento real de la empresa por arriba de la inflacion. & c) Tasa Minima Aceptable de Rendimiento (TMAR) & Es la TMAR, referencia minima de rendimiento aceptable. \\
51 & Plan al que se le asigna capital e insumos materiales, humanos y tecnicos. Su objetivo es generar un rendimiento economico a un determinado plazo. & d) Proyecto de inversion & Es un proyecto de inversion. \\
52 & Las variables o factores conforman un campo dinamico e intenso de fuerzas que asumiendo tendencias y direcciones afectan a las organizaciones. & Verdadero & El ambiente externo influye en decisiones y resultados. \\
53 & Los objetivos empresariales son: Economicos, sociales y tecnicos. & Verdadero & Es una clasificacion tradicional de objetivos de empresa. \\
54 & En la practica, el nombre y el numero de las areas funcionales cambian segun la organizacion, el tamano de la empresa; aunque la estructura y los nombres cambian. & Verdadero & La estructura se adapta al giro, tamano y necesidades de cada empresa. \\
\end{longtable}
\endgroup
\end{landscape}
\clearpage
\pagestyle{fancy}

\subsection{Analisis del cuestionario}

El desarrollo muestra que el cuestionario sintetiza varias capas de formacion administrativa. No solo revisa conceptos introductorios, sino que obliga a enlazar planeacion, control, comercializacion, analisis financiero y evaluacion de inversiones bajo un mismo criterio de gestion. Esa integracion es relevante porque en la practica directiva los problemas no llegan separados por asignatura: llegan como decisiones que exigen comprender mercado, recursos, costos, rentabilidad y ejecucion al mismo tiempo \citep{amelicaProcesoAdmin2019,snhuAdminMerca2023}.

Desde esta lectura, las preguntas numericas cumplen una funcion estrategica: convierten nociones abstractas en evidencia verificable. Calcular margen neto, valor presente, razon circulante o costo de capital obliga a pasar del lenguaje conceptual al criterio cuantitativo. En consecuencia, el producto demuestra que la formacion administrativa requiere precision tecnica, pero tambien capacidad para interpretar que significa cada resultado dentro de una organizacion.

\section{Conclusion}

La actividad confirma que los fundamentos de administracion, finanzas, mercadotecnia y proyectos de inversion no son bloques aislados, sino componentes de una misma logica de gestion. Resolver el cuestionario con justificacion permite identificar que la planeacion orienta la accion, el control mide desviaciones, la mercadotecnia traduce valor hacia el mercado y las finanzas verifican la viabilidad de las decisiones. Ese aprendizaje tiene implicaciones directas para la formacion en la maestria, porque fortalece la capacidad de argumentar, calcular y decidir con mayor consistencia. En terminos de criterio final, el valor academico del ejercicio no esta en acumular aciertos, sino en demostrar que cada respuesta puede sostenerse con comprension tecnica y sentido estrategico.

\clearpage
\nocite{itescaSitioWeb}
\bibliography{\itescabibliography}

\end{document}


